\chapter{Lezione 18}
\label{chap:lezione_18} % Etichetta per riferirsi al capitolo

\begin{flushright}
\textit{Data: 12/05/2025}
\end{flushright}
\textit{Bibliografia: 
Lagevin equation as an effective description. Examples from harmonic trap and bistable potentials. Fully-connected Ising model, Landau Theory, and Langevin dynamics of the magnetization in the homogeneous case}

\section*{Recap delle lezioni precedenti e fondamenti}

In precedenza, ci siamo concentrati sulla Master Equation per sistemi discreti. Tuttavia, la forma differenziale dell'equazione di Chapman-Kolmogorov che avevamo ottenuto è più generale e potente. Ricordiamo brevemente che l'equazione generale ha una forma che include sia i processi discontinui/di salto (master equation) e sia l'espansione di Kramers-Moyal per i processi continui.

Se consideriamo solo processi continui, la parte relativa ai salti viene eliminata. L'equazione di Kramers-Moyal descrive l'evoluzione della densità di probabilità $P(x,t)$. Una forma generale può essere scritta considerando i momenti del processo. I coefficienti $a^{(m)}(x,t)$ sono definiti come i momenti locali del tasso di variazione della variabile stocastica $x$:
\begin{equation}
a^{(m)}(x,t) \equiv \frac{1;}{m!} \lim_{\Delta t\rightarrow0} \frac{1}{\Delta t} \lim_{\epsilon\rightarrow0} \int_{|y-x|\le\epsilon} (y-x)^{m} P(y,t+\Delta t|x,t) \, dy
\end{equation}
Questi $a^{(m)}$ sono i momenti del rate della nostra distribuzione.

L'equazione di Kramers-Moyal può essere scritta come:
\begin{equation}
\frac{\partial P(x,t)}{\partial t} = \sum_{m=1}^{\infty} (-1)^{m} \frac{\partial^m}{\partial x^m} [a^{(m)}(x)P(x,t)]
\end{equation}

L'equazione di Fokker-Planck è un caso particolare dell'equazione di Kramers-Moyal, ottenuta troncando l'espansione al secondo termo. Questo è valido quando i coefficienti $a^{(m)}$ sono nulli per $m > 2$.

Questo si verifica, ad esempio, per processi guidati da rumore Gaussiano bianco, come nell'equazione di Langevin. Per un processo di Langevin con rumore Gaussiano, i momenti scalano come $\langle(\Delta x)^{n}\rangle$. Più specificamente, per i momenti pari $\langle(\Delta x)^{2n}\rangle \sim (\Delta t)^{n}$.

Considerando il limite per $\Delta t \rightarrow 0$:
\begin{equation}
\lim_{\Delta t\rightarrow0}\frac{1}{\Delta t}\langle(\Delta x)^{2n}\rangle =
\begin{cases}
\text{costante} \neq 0, & \text{se } n=1 \\
0, & \text{se } n>1
\end{cases}
\end{equation}

Questo implica che solo $a^{(1)}$ (drift) e $a^{(2)}$ (diffusione) sono non nulli nel limite diffusivo, mentre i termini superiori si annullano. Di conseguenza, l'equazione di Fokker-Planck è:
\begin{equation}
\partial_{t}P(x,t)=-\partial_{x}(a^{(1)}(x)P(x,t))+\partial_{x}^{2}(a^{(2)}(x)P(x,t))
\end{equation}

Per sistemi discreti, avevamo introdotto il teorema H utilizzando la divergenza di Kullback-Leibler:
\begin{equation}
H_{KL}(t) \equiv \sum_{n} p_{n}(t) \log\frac{p_{n}(t)}{p_{n}^{\infty}}
\end{equation}
dove $p_n^{\infty}$ è la distribuzione stazionaria. Si ha che $\frac{dH_{KL}}{dt} \le 0$, e per $t\rightarrow\infty$, $p_{n}(t)\rightarrow p_{n}^{\infty}$ e $H_{KL} \rightarrow 0$.

Si può anche definire l'entropia di Shannon:
\begin{equation}
H(t) \equiv -\sum_{n} P_{n}(t) \log P_{n}(t)
\end{equation}
L'entropia tende a crescere per sistemi isolati che evolvono verso l'equilibrio. La funzione H di Boltzmann, o la divergenza di Kullback-Leibler, misura la ``distanza'' dalla distribuzione di equilibrio e decresce nel tempo.

\section{Operatore di Fokker-Planck}
Possiamo definire l'operatore di Fokker-Planck $\mathcal{L}_{\text{FP}}$ tale che:
\begin{equation}
\partial_{t}P(x,t) = \mathcal{L}_{\text{FP}}P(x,t)
\end{equation}

L'operatore è dato da:
\begin{equation}
\mathcal{L}_{\text{FP}} \equiv -\frac{\partial}{\partial x}a^{(1)}(x) + \frac{\partial^{2}}{\partial x^{2}}a^{(2)}(x)
\end{equation}

Una distribuzione stazionaria $P_{\infty}(x)$ è una soluzione dell'equazione di Fokker-Planck che non dipende dal tempo, cioè $\partial_t P_{\infty}(x) = 0$. Quindi:
\begin{equation}
\mathcal{L}_{\text{FP}}P_{\infty}(x)=0
\end{equation}

\begin{tcolorbox}[
    colback=colorD!5,      
    colframe=colorD,       
    title={Notazione},
    fonttitle=\bfseries\sffamily,
    arc=3mm,               
    boxrule=1.2pt,         
    halign title=flush left
]
La distribuzione stazionaria $P_{\infty}(x) = \lim_{t\rightarrow\infty}P(x,t)$ è l'autofunzione dell'operatore di Fokker-Planck corrispondente a un autovalore nullo.
\end{tcolorbox}

L'operatore di Fokker-Planck $\mathcal{L}_{\text{FP}}$ in generale non è simmetrico (hermitiano). Questo complica l'analisi spettrale rispetto agli operatori simmetrici.

La soluzione generale $P(x,t)$ può essere espressa come uno sviluppo in autofunzioni $\psi_m(x)$ dell'operatore $\mathcal{L}_{\text{FP}}$:
\begin{equation}
P(x,t)=\sum_{m} c_m \psi_{m}(x)e^{-\lambda_{m}t}
\end{equation}
dove $c_m$ sono coefficienti determinati dalle condizioni iniziali, e le $\psi_m(x)$ e $\lambda_m$ soddisfano l'equazione agli autovalori:
\begin{equation}
\mathcal{L}_{\text{FP}}\psi_{m}(x)=-\lambda_{m}\psi_{m}(x)
\end{equation}

Gli autovalori $\lambda_m$ sono i tassi di rilassamento inversi:
\begin{equation} 
\lambda_m = \frac{1}{\tau_m} 
\end{equation}

\begin{tcolorbox}[
    colback=colorD!5,      
    colframe=colorD,       
    title={Importante},
    fonttitle=\bfseries\sffamily,
    arc=3mm,               
    boxrule=1.2pt,         
    halign title=flush left
]
Affinché esista una distribuzione stazionaria non banale ($P_\infty(x) \neq 0$) a tempi lunghi, l'autovalore più piccolo (associato al tempo di rilassamento più lungo) deve essere nullo, cioè $\lambda_0=0$.
Se tutti gli autovalori fossero strettamente positivi, ogni termine $e^{-\lambda_m t}$ tenderebbe a zero per $t\rightarrow\infty$, e quindi $P(x,t)\rightarrow 0$.
L'autofunzione $\psi_0(x)$ associata a $\lambda_0=0$ è la distribuzione stazionaria $P_\infty(x)$.
\end{tcolorbox}

\subsection{Caso con Potenziale e Operatore Hermitiano Associato}
Una situazione di particolare interesse fisico si ha quando il coefficiente di drift $a^{(1)}(x)$ deriva da un potenziale $\Phi(x)$:
\begin{equation}
a^{(1)}(x)=-\frac{\partial\Phi(x)}{\partial x}
\end{equation}
e il coefficiente di diffusione $a^{(2)}(x)=D$ è costante.

In questo caso, la soluzione stazionaria dell'equazione di Fokker-Planck è data dalla distribuzione di Boltzmann:
\begin{equation}
P_{\infty}(x)=N e^{-\Phi(x)/D}
\end{equation}
dove $N$ è una costante di normalizzazione. Per la normalizzabilità, si richiede che il potenziale diverga sufficientemente rapido all'infinito:
\begin{equation}
\lim_{x\rightarrow\pm\infty}\Phi(x)=+\infty
\end{equation}

Anche se $\mathcal{L}_{\text{FP}}$ non è simmetrico, in questo caso è possibile associargli un operatore hermitiano $\mathcal{H}$ tramite una trasformazione di similarità:
\begin{equation}
\mathcal{H} = e^{\Phi(x)/(2D)} \mathcal{L}_{\text{FP}} e^{-\Phi(x)/(2D)}
\end{equation}

Le autofunzioni $\varphi_n(x)$ di $\mathcal{H}$ sono legate alle autofunzioni $\psi_n(x)$ di $\mathcal{L}_{\text{FP}}$ da:
\begin{equation}
\varphi_{n}(x)=e^{\Phi(x)/(2D)} \psi_{n}(x)
\end{equation}

L'operatore $\mathcal{H}$ e $\mathcal{L}_{\text{FP}}$ hanno lo stesso spettro di autovalori $\{\lambda_n\}$.

\begin{tcolorbox}[
    colback=colorD!5,      
    colframe=colorD,       
    title={Proprietà},
    fonttitle=\bfseries\sffamily,
    arc=3mm,               
    boxrule=1.2pt,         
    halign title=flush left
]
Poiché $\mathcal{H}$ è hermitiano, si può dimostrare che l'autovalore più piccolo $\lambda_0=0$ è unico, garantendo l'unicità della soluzione stazionaria $P_\infty(x)$.
\end{tcolorbox}

\section{Equazione di Langevin Generalizzata}
L'equazione di Langevin generalizzata per una variabile $x(t)$ è:
\begin{equation}
\dot{x}(t) = f(x(t)) + \eta(t)
\end{equation}
dove $f(x)$ rappresenta le forze sistematiche (deterministiche) e $\eta(t)$ è un termine di rumore stocastico. Tipicamente, per il rumore si assume:
\begin{itemize}
    \item Media nulla: $\langle\eta(t)\rangle=0$ (eventuali medie non nulle possono essere riassorbite in $f(x)$).
    \item Correlazione delta nel tempo (rumore bianco): $\langle\eta(t)\eta(s)\rangle=2D\delta(t-s)$, dove $D$ è l'intensità del rumore (coefficiente di diffusione).
\end{itemize}

Questa descrizione è valida quando esiste una separazione di scale temporali: il tempo di correlazione del rumore $\tau_R$ è molto più piccolo del tempo caratteristico della dinamica di $x$, $\tau_D$.

Questo tipo di dinamica è chiamata dinamica di rilassamento, poiché il sistema tende a rilassare verso stati di equilibrio o stati stazionari sotto l'influenza delle forze e del rumore.

\subsection{Moto Browniano Libero ($f(x)=0$)}
Se $f(x)=0$, l'equazione diventa $\dot{x}=\eta(t)$, che descrive il moto Browniano standard (processo di Wiener se $\eta(t)$ è rumore bianco Gaussiano).

\subsection{Trappola Armonica}
Consideriamo una particella in un potenziale armonico $H(x)=\frac{1}{2}kx^{2}$. La forza deterministica è $f(x) = -H'(x) = -kx$. L'equazione di Langevin è:
\begin{equation}
\dot{x}=-kx+\eta(t)
\end{equation}
con una condizione iniziale $x(t_0)=x_0$.

\paragraph{Caso senza rumore ($\eta(t)=0$):}
L'equazione diventa $\dot{x}=-kx$. Il punto di equilibrio si ha per $H'(x)=0$, cioè $x_{\infty}=0$. La soluzione è $x(t)=x_{0}e^{-k(t-t_0)}$. Per $t\rightarrow\infty$, $x(t)\rightarrow 0$. La distribuzione di probabilità stazionaria è una delta di Dirac centrata in zero: $P_{\infty}(x)=\delta(x)$.

\paragraph{Caso con rumore ($\eta(t)\neq 0$):}
La soluzione dell'equazione di Langevin lineare è:
\begin{equation}
x(t)=x_{0}e^{-k(t-t_0)}+\int_{t_0}^{t}ds\, e^{-k(t-s)}\eta(s)
\label{eq:sol_armonica}
\end{equation}

Il valore medio di $x(t)$ evolve come: $\langle x(t)\rangle = x_0 e^{-k(t-t_0)} + \int_{t_0}^{t}ds\, e^{-k(t-s)}\langle\eta(s)\rangle$. Poiché $\langle\eta(s)\rangle=0$:
\begin{equation}
\langle x(t)\rangle = x_0 e^{-k(t-t_0)}
\end{equation}
Per $t\rightarrow\infty$, $\langle x(t)\rangle \rightarrow 0$.

La distribuzione stazionaria è una Gaussiana centrata in zero, con varianza $\sigma^2 = D/k$:
\begin{equation}
P_{\infty}(x) \propto e^{-\frac{kx^{2}}{2D}} \propto e^{-\frac{H(x)}{D}}
\end{equation}
Questo mostra che anche in presenza di rumore, il sistema tende a fluttuare attorno al minimo del potenziale. La larghezza della distribuzione dipende dall'intensità del rumore $D$ e dalla ``rigidità'' della trappola $k$.

\subsection{Potenziale Bistabile}
Consideriamo un potenziale a doppia buca (bistabile):
\begin{equation}
H(x)=-\frac{a}{2}x^{2}+\frac{b}{4}x^{4}
\end{equation}
con $a,b >0$. Questo potenziale ha due minimi stabili e un massimo locale instabile. L'equazione di Langevin è $\dot{x}=-H'(x)+\eta(t)$.

\paragraph{Analisi senza rumore ($\eta(t)=0$):}
I punti di equilibrio si trovano da $H'(x)=0$:
\begin{equation}
H'(x)=-ax+bx^{3}=x(-a+bx^{2})=0
\end{equation}
Le soluzioni sono: $x_1=0$ e $x_{2,3}=\pm\sqrt{\frac{a}{b}}$.

Per studiare la stabilità, calcoliamo la derivata seconda $H''(x)=-a+3bx^{2}$:
\begin{itemize}
    \item $H''(0)=-a < 0$ : $x_1$ è un punto di equilibrio instabile (massimo locale).
    \item $H''(x_{2,3})=-a+3b\left(\frac{a}{b}\right)=2a > 0$ : $x_{2,3}$ sono punti di equilibrio stabili (minimi locali).
\end{itemize}
Una particella posta in uno dei minimi, senza rumore, vi rimarrebbe indefinitamente. Se perturbata leggermente, scivolerebbe verso il fondo della buca.

\paragraph{Caso con rumore ($D \ne 0$):}

In presenza di rumore, la particella può saltare da una buca di potenziale all'altra. Questo fenomeno è spesso mediato da fluttuazioni sufficientemente grandi da superare la barriera di potenziale tra i minimi. Tali eventi rari ma importanti sono chiamati istantoni.


L'altezza della barriera di potenziale, misurata da un minimo $x_m = \pm\sqrt{a/b}$ al massimo locale $x=0$, è: $\Delta E_b = H(0) - H(x_m)$.
\begin{equation*}
H(x_m) = H\left(\pm\sqrt{\frac{a}{b}}\right) = -\frac{a}{2}\left(\frac{a}{b}\right)+\frac{b}{4}\left(\frac{a^2}{b^2}\right) = -\frac{a^2}{2b}+\frac{a^2}{4b} = -\frac{a^2}{4b}
\end{equation*}
\begin{equation*}
H(0) = 0
\end{equation*}

Quindi, la barriera che la particella deve superare per andare da un minimo all'altro (passando per $x=0$) è:
\begin{equation}
\Delta E_b = H(0) - H\left(\pm\sqrt{\frac{a}{b}}\right) = 0 - \left(-\frac{a^2}{4b}\right) = \frac{a^2}{4b}
\end{equation}

Il tempo medio di salto da una buca all'altra, $\tau_{\text{jump}}$, è descritto dalla legge di Arrhenius:
\begin{equation}
\tau_{\text{jump}} \propto e^{\frac{\Delta E_b}{D}}
\end{equation}
Questo indica che il tempo di salto cresce esponenzialmente con l'altezza della barriera e decresce con l'aumentare dell'intensità del rumore $D$.

La distribuzione stazionaria per questo sistema è ancora data da $P_\infty(x) \propto e^{-H(x)/D}$, che avrà due picchi corrispondenti ai minimi del potenziale. 

Un potenziale della forma $H(x)=\alpha\frac{x^{2}}{2}+\beta\frac{x^{4}}{4}$ (con $\beta>0$) può descrivere sia una singola buca (se $\alpha>0$) sia una doppia buca (se $\alpha<0$, ponendo $\alpha = -a$). Questi sistemi bistabili sono prototipi per molti fenomeni fisici, incluse le transizioni di fase.

\section{Modello di Landau}
Il potenziale a doppia buca emerge naturalmente nella teoria di Landau delle transizioni di fase. Questa teoria descrive il comportamento di un sistema vicino a un punto critico in termini di un parametro d'ordine $\varphi$.

Consideriamo il modello di Ising come esempio. L'Hamiltoniana è:
\begin{equation}
\mathcal{H}[S]=-\sum_{i,j}J_{ij}S_{i}S_{j}-h\sum_{i}S_{i}
\end{equation}
dove $S_i=\pm 1$ sono gli spin, $J_{ij}$ l'accoppiamento tra spin $i$ e $j$, e $h$ un campo magnetico esterno.
\begin{equation} 
J_{ij} = \begin{cases} J \lessgtr 0 \quad & (i,j) \text{ primi vicini} \\ 0 \quad &\text{altrimenti} \end{cases} 
\end{equation}
Se $J_{ij}>0$ il modello descrive un ferromagnete, altrimenti un antiferromagnete.

Il modello di Ising 1D (con interazione solo tra primi vicini) non presenta una transizione di fase a temperatura finita, cioè $\langle S_i \rangle = 0$ per $h=0, T>0$. In dimensioni superiori ($D \ge 2$), o nel limite di campo medio ($D\rightarrow\infty$ o modello fully-connected), possono verificarsi transizioni di fase.

Nel modello fully-connected (o di Curie-Weiss), tutti gli spin interagiscono con tutti gli altri con la stessa forza, riscalata per il numero di spin $N$ per garantire un limite termodinamico sensato. Per $h=0$:
\begin{equation}
\mathcal{H}[S]=-\frac{J}{2N}\sum_{i,j}S_{i}S_{j}
\end{equation}

L'energia libera per spin nel limite termodinamico ($N\rightarrow\infty$) è:
\begin{equation} 
f(\beta) = -\frac{1}{N\beta}\ln Z_{\beta} 
\end{equation}
dove $Z_\beta$ è la funzione di partizione:
\begin{tcolorbox}[
    colback=colorA!5,      
    colframe=colorA,       
    arc=3mm,               
    boxrule=1.2pt
]
\begin{equation}
Z_{\beta}=\text{Tr}_{\{S\}}e^{-\beta \mathcal{H}} = \text{Tr}_{\{S\}}e^{\frac{\beta J}{2N}\sum_{i,j}S_{i}S_{j}}
\end{equation}
\end{tcolorbox}
dove $\text{Tr}_{\{S\}} \equiv \sum_{S_1=\pm1} \dots \sum_{S_N=\pm1}$.

Usando l'identità $\sum_{i,j}S_{i}S_{j} = (\sum_{i}S_{i})^{2} - \sum_{i}S_{i}^{2} = (\sum_{i}S_{i})^{2}-N$ (poiché $S_i^2=1$):
\begin{equation}
Z_{\beta} = e^{-\frac{\beta J}{2}} \text{Tr}_{\{S\}} \exp\left(\frac{\beta J}{2N}\left(\sum_{i}S_{i}\right)^{2}\right)
\end{equation}

Definiamo il parametro d'ordine, ovvero la magnetizzazione per spin:
\begin{equation} 
\varphi = \frac{1}{N}\sum_{i}S_{i} 
\end{equation}

Per calcolare $Z_\beta$, si utilizza la trasformazione di Hubbard-Stratonovich, introducendo una variabile ausiliaria $\hat{\varphi}$ e una delta di Dirac (il termo $e^{-\beta J/2}$ si riassorbe nelle costanti):
\begin{equation}
Z_{\beta} = \text{Tr}_{\{S\}} \int d \varphi \; \delta\left(N\varphi-\sum_{i} S_{i}\right) \; e^{\frac{\beta J}{2N}(N \varphi)^{2}}
\end{equation}

Si scrive la delta come integrale:
\begin{equation} 
\delta\left(N\varphi-\sum_{i} S_{i}\right)=\int_{-i\infty}^{+i\infty}\frac{d\hat{\varphi}}{2\pi i}e^{-\hat{\varphi}(N\varphi-\sum_{i}S_i)} 
\end{equation}

\begin{equation}
Z_{\beta} = \text{Tr}_{\{S\}} \int d \varphi \; \int_{-i\infty}^{+i\infty}\frac{d\hat{\varphi}}{2\pi i}e^{-\hat{\varphi}(N\varphi-\sum_{i}S_i)} \; e^{\frac{\beta J}{2N}(N \varphi)^{2}}
\end{equation}

Si esegue la traccia sugli spin:
\begin{align} 
\text{Tr}_{\{S\}}e^{\hat{\varphi}\sum_{i} S_i} &= \sum_{S_1=\pm1} \dots \sum_{S_N=\pm1} e^{\hat{\varphi}(S_1 + \dots + S_N)} \nonumber \\ 
&= \Bigl( \sum_{S=\pm1} e^{\hat{\varphi} S} \Bigr)^N = \Bigl( e^{\hat{\varphi}} + e^{-\hat{\varphi}} \Bigr)^N \nonumber \\
&=(2\cosh\hat{\varphi})^{N} = \exp \Bigl(N \ln(2 \cosh\hat{\varphi}) \Bigr)
\end{align}

La funzione di partizione diventa:
\begin{equation}
Z_{\beta} \approx C \int d\varphi \; d\hat{\varphi} \; e^{-NG(\varphi,\hat{\varphi})}
\label{eq:landau_40}
\end{equation}
dove:

\begin{align}
G(\varphi,\hat{\varphi}) &= -\frac{\beta J}{2}\varphi^{2}+\hat{\varphi}\varphi-\ln(2\cosh\hat{\varphi}) \nonumber \\ 
&= -\frac{\beta J}{2}\varphi^{2}+\hat{\varphi}\varphi-\ln(2) -\ln(\cosh\hat{\varphi}) \nonumber \\ 
&\approx -\frac{\beta J}{2}\varphi^{2}+\hat{\varphi}\varphi -\ln(\cosh\hat{\varphi})
\label{eq:landau_41}
\end{align}
Nel passaggio dalla (\ref{eq:landau_40}) alla (\ref{eq:landau_41}) abbiamo ritenuto trascurabile il termine $\ln(2)$.

Nel limite $N\rightarrow\infty$, l'integrale è dominato dal metodo del punto di sella:
\begin{equation}
Z_{\beta} \approx C \; e^{-NG_{\text{SP}}}
\end{equation}
dove $G_{\text{SP}}$ è $G$ valutata al punto di sella.

L'energia libera per spin è:
\begin{equation} 
f(\beta) = -\frac{1}{N\beta}\ln Z_{\beta} \approx \frac{G_{\text{SP}}}{\beta} 
\end{equation}

Le equazioni del punto di sella sono:
\begin{equation} 
\left. \frac{\partial G}{\partial\varphi}\right|_{\text{SP}}=0 \quad \longrightarrow \quad \hat{\varphi}_{\text{SP}}=\beta J \varphi_{\text{SP}} 
\end{equation}
\begin{equation} 
\left. \frac{\partial G}{\partial\hat{\varphi}}\right|_{\text{SP}}=0 \quad \longrightarrow \quad \varphi_{\text{SP}} = \tanh(\hat{\varphi}_{\text{SP}}) 
\end{equation}

Combinandole, si ottiene l'equazione di autoconsistenza per la magnetizzazione $\varphi_{\text{SP}}$:
\begin{tcolorbox}[
    colback=colorA!5,      
    colframe=colorA,       
    arc=3mm,               
    boxrule=1.2pt
]
\begin{equation}
\varphi_{\text{SP}} = \tanh(\beta J\varphi_{\text{SP}})
\end{equation}
\end{tcolorbox}
Questa prende il nome di \textbf{Equazione di Curie-Weiss}.

Il sistema descritto da questa equazione di autoconsistenza va incontro a una transizione di fase del second'ordine. Le transizioni del second'ordine sono caratterizzate dal fatto che il parametro d'ordine, in questo caso la magnetizzazione $\varphi$, varia con continuità attraverso la temperatura critica $T_c$. Nello specifico:
\begin{itemize}
    \item Per temperature $T > T_c$ (fase disordinata o simmetrica), il parametro d'ordine è nullo: $\varphi = 0$. Questa fase possiede una simmetria più elevata.
    \item Per temperature $T < T_c$ (fase ordinata o a simmetria rotta), il parametro d'ordine assume un valore non nullo: $\varphi \neq 0$. La comparsa di un $\varphi \neq 0$ segnala una rottura spontanea della simmetria presente nella fase ad alta temperatura. Ad esempio, il sistema sceglie una particolare direzione di magnetizzazione.

\item Alla temperatura critica $T = T_c$, il parametro d'ordine è ancora nullo, ma inizia a deviare da zero per $T$ immediatamente inferiore a $T_c$.
\end{itemize}

L'idea fondamentale di Landau fu di assumere che, poiché $\varphi$ è continuo e piccolo vicino a $T_c$, l'energia libera del sistema potesse essere sviluppata in serie di potenze di $\varphi$.

\section{Sviluppo dell'Energia Libera di Landau}
L'energia libera di Landau (o potenziale di Landau-Ginzburg) si ottiene sostituendo $\hat{\varphi}_{\text{SP}}$ in $G(\varphi, \hat{\varphi})$ e sviluppando per $\varphi$ piccolo, vicino alla temperatura critica $T_c$.

Sostituendo la relazione dell'equazione di sella per il campo ausiliario nell'equazione per $G_{\text{SP}}$ si ha:
\begin{equation}
G_{\text{SP}}(\varphi_{\text{SP}}) = \frac{\beta J}{2}\varphi_{\text{SP}}^{2}-\ln \Bigl(\cosh(\beta J \varphi_{\text{SP}}) \Bigr)
\end{equation}

L'energia libera per spin è data da:
\begin{equation}
f(\beta,\varphi) = \frac{1}{\beta} G_{\text{SP}}(\varphi) = \frac{J}{2}\varphi_{\text{SP}}^{2}-\frac{1}{\beta} \ln \Bigl(\cosh(\beta J \varphi_{\text{SP}}) \Bigr)
\end{equation}
dove il primo termine è di tipo energetico, mentre il secondo è di tipo entropico.

Per $\varphi$ piccolo, usiamo il seguente sviluppo in serie di Taylor:
\begin{equation} 
\ln(\cosh x) \approx \frac{x^2}{2} - \frac{x^4}{12} + O(x^6) 
\end{equation}

L'energia libera diventa quindi:
\begin{align}
f(\beta,\varphi) &\approx \frac{J}{2}\varphi_{\text{SP}}^{2}-\frac{1}{\beta} \Bigl[ \frac{(\beta J \varphi_{\text{SP}})^2 }{2}-\frac{(\beta J \varphi_{\text{SP}})^4 }{12}\Bigr] \nonumber \\ 
&= \frac{J}{2}(1 - \beta J)\varphi_{\text{SP}}^{2} + \frac{J}{12} (J \beta)^3 \varphi_{\text{SP}}^{4}
\end{align}

L'energia libera di Landau è tipicamente scritta come:
\begin{tcolorbox}[
    colback=colorA!5,      
    colframe=colorA,       
    arc=3mm,               
    boxrule=1.2pt
]
\begin{equation}
f(\varphi)\equiv H_{LG}(\varphi) = \frac{a(T)}{2}\varphi^{2}+\frac{b}{4}\varphi^{4}
\label{eq:landau_ginzburg_pot}
\end{equation}
\end{tcolorbox}

Confrontando i termini si ricavano i coefficienti:
\begin{tcolorbox}[
    colback=colorA!5,      
    colframe=colorA,       
    arc=3mm,               
    boxrule=1.2pt
]
\begin{align}
a(T) &= J(1-\beta J) \\
b &= \frac{J^4 \beta^3}{3}
\end{align}
\end{tcolorbox}

La temperatura critica $T_c$ è definita dalla condizione $a(T_c)=0$, da cui si ottiene $k_B T_c = J$.

Il coefficiente $a(T)$ cambia segno in corrispondenza di $T_c$:
\begin{itemize}
    \item $a(T)>0$ per $T>T_c$ (fase disordinata, $\varphi=0$)
    \item $a(T)<0$ per $T<T_c$ (fase ordinata, $\varphi\neq0$)
\end{itemize}
Il coefficiente $b$ è strettamente positivo ($b>0$) vicino a $T_c$, garantendo la stabilità del sistema per grandi valori di $\varphi$. 
\noindent Questo funzionale $H_{LG}(\varphi)$ assume la stessa forma del potenziale bistabile discusso in precedenza quando $a(T)<0$.

\subsection{Analisi dell'Energia Libera di Landau}
I valori di equilibrio del parametro d'ordine $\bar{\varphi}$ minimizzano il potenziale $H_{LG}(\varphi)$:
\begin{equation} 
\left. \frac{\partial H_{LG}}{\partial\varphi}\right|_{\bar{\varphi}} = a(T)\bar{\varphi}+b\bar{\varphi}^{3}=0 
\end{equation}

Le soluzioni formali sono:
\begin{itemize}
    \item $\bar{\varphi}_{0}=0$
    \item $\bar{\varphi}_{1,2}=\pm\sqrt{\frac{-a(T)}{b}}$ (reali e non nulle solo se $a(T)<0$, ovvero per $T<T_c$)
\end{itemize}

La stabilità locale è determinata dal segno della derivata seconda $H_{LG}'' = a(T)+3b\bar{\varphi}^{2}$:
\begin{itemize}
    \item \textbf{Caso} $T>T_c \implies a(T)>0$:\\
    L'unica soluzione reale è $\bar{\varphi}_0=0$. Si ha $H_{LG}''(0)=a(T)>0$, quindi $\bar{\varphi}_0=0$ è un minimo stabile. Il potenziale presenta una singola buca.
    \item \textbf{Caso} $T=T_c \implies a(T)=0$:\\
    Il potenziale si riduce a $H_{LG}(\varphi) = \frac{b}{4}\varphi^4$. L'unica soluzione è $\bar{\varphi}_0=0$ con $H_{LG}''(0)=0$. Il potenziale è estremamente piatto nell'intorno dello zero.
    \item \textbf{Caso} $T<T_c \implies a(T)<0$:\\
    La soluzione $\bar{\varphi}_0=0$ comporta $H_{LG}''(0)=a(T)<0$, qualificandosi come un massimo locale instabile. Le soluzioni coerenti $\bar{\varphi}_{1,2}= \pm\sqrt{\frac{|a(T)|}{b}}$ restituiscono:
    \begin{equation}
    H_{LG}''(\bar{\varphi}_{1,2})=a(T)+3b\left(\frac{-a(T)}{b}\right)=-2a(T)=2|a(T)|>0
    \end{equation}
    Questi punti rappresentano due minimi stabili degeneri. Il potenziale assume una forma a doppia buca.
\end{itemize}

\subsection{Inclusione del Campo Magnetico ed Esponenti Critici}
Se si include l'effetto di un campo magnetico esterno $h$, l'energia libera di Landau viene modificata dall'aggiunta di un termine di accoppiamento lineare:
\begin{equation}
H_{LG}(\varphi) = \frac{a(T)}{2}\varphi^{2}+\frac{b}{4}\varphi^{4}-h\varphi
\end{equation}

L'equazione di stato si ricava imponendo l'estremizzazione del funzionale:
\begin{equation} 
\frac{\partial H_{LG}}{\partial\varphi}=a(T)\varphi+b\varphi^{3}-h=0 
\label{eq:stato_landau}
\end{equation}

Sull'isoterma critica ($T=T_c \implies a(T)=0$), l'equazione (\ref{eq:stato_landau}) si riduce a $b\bar{\varphi}^3=h$, da cui consegue che $\bar{\varphi} \sim h^{1/3}$. L'esponente critico di campo associato è $\delta_{MF}=3$.

La suscettività magnetica $\chi$ si ottiene differenziando l'equazione di stato:
\begin{equation} 
\chi^{-1} = \left(\frac{\partial \varphi}{\partial h}\right)^{-1}=a(T)+3b\bar{\varphi}^2
\end{equation}

\begin{itemize}
    \item Per $T>T_c$, dato che $\bar{\varphi}=0$, si ha $\chi^{-1}=a(T)$. Poiché l'espansione lineare prevede $a(T) \propto (T-T_c)$, la suscettività diverge come $\chi \propto (T-T_c)^{-\gamma}$ con un esponente critico $\gamma_{MF}=1$.
    \item Per $T<T_c$, sostituendo il valore di equilibrio del parametro d'ordine si ottiene $\chi^{-1}=a(T)+3b(-a(T)/b) = -2a(T)$, confermando la divergenza della suscettività con la medesima legge di potenza e lo stesso esponente critico $\gamma_{MF}=1$.
\end{itemize}

La dinamica del parametro d'ordine $\varphi(t)$ in prossimità del punto critico può essere descritta da un'equazione di Langevin strutturata direttamente sul potenziale energetico di Landau:
\begin{equation}
\frac{d\varphi}{dt} = -\frac{\partial H_{LG}}{\partial \varphi} + \eta_{\varphi}(t)
\end{equation}
dove $\eta_{\varphi}(t)$ rappresenta il termine di rumore termico gaussiano e bianco correlato.

Questa equazione differenziale stocastica descrive il rilassamento del parametro d'ordine (inteso come grado di libertà lento del sistema collettivo) verso il suo valore di equilibrio macroscopico, fluttuando sotto l'azione del bagno termico. Il principio di universalità assicura che questo formalismo dinamico si applichi a una vasta classe di sistemi fisici differenti durante transizioni di fase continue.